% Môn: Vật lý
% Lớp: 10
% Chương: Chương I. Mở đầu
% Bài: L10C1 Bài 1. Làm quen với Vật lí
% Loại: Lý thuyết
% File riêng, không trộn vào de.tex — thay nội dung khi soạn xong
% Định nghĩa lệnh Lý thuyết — dùng khi biên dịch lt.tex bằng TeX.
% App web đọc lt.tex trực tiếp (bỏ \input file này).
% Trong lt.tex, từ thư mục bài:
%   % Định nghĩa lệnh Lý thuyết — dùng khi biên dịch lt.tex bằng TeX.
% App web đọc lt.tex trực tiếp (bỏ \input file này).
% Trong lt.tex, từ thư mục bài:
%   % Định nghĩa lệnh Lý thuyết — dùng khi biên dịch lt.tex bằng TeX.
% App web đọc lt.tex trực tiếp (bỏ \input file này).
% Trong lt.tex, từ thư mục bài:
%   \input{../../../../_lenh/lythuyet.tex}
% từ gốc repo:
%   \input{ngan-hang/_lenh/lythuyet.tex}
\makeatletter
\providecommand{\indam}[1]{\textbf{#1}}
\providecommand{\chuy}[1]{\par\smallskip\noindent\textit{#1}\par}
\providecommand{\luuy}[1]{\par\smallskip\noindent\textbf{Lưu ý.} #1\par}
\providecommand{\ghichu}[1]{\par\smallskip\noindent\textbf{Ghi chú.} #1\par}
\providecommand{\vidu}[1]{\par\smallskip\noindent\textbf{Ví dụ.} #1\par}
\providecommand{\Blythuyet}{\subsection*{Lý thuyết}}
% Hai cột: kiến thức | hình
\providecommand{\haicotchay}[2]{%
  \par\medskip\noindent
  \begin{minipage}[t]{0.52\linewidth}#1\end{minipage}\hfill
  \begin{minipage}[t]{0.46\linewidth}#2\end{minipage}%
  \par\medskip
}
\providecommand{\kienthuccot}[2]{\haicotchay{#1}{#2}}
% Dự phòng nếu chưa nạp graphicx / caption (không ghi đè bản thật)
\providecommand{\resizebox}[3]{#3}
\providecommand{\captionof}[2]{\par\smallskip\noindent\textit{#2}\par}
\newenvironment{hoatdong}[1]{%
  \par\medskip\noindent\textbf{Hoạt động.} #1\par\smallskip
}{\par\medskip}
\newenvironment{emcobiet}{%
  \par\medskip\noindent\textbf{Em có biết}\par\smallskip
}{\par\medskip}
\newenvironment{emdahoc}{%
  \par\medskip\noindent\textbf{Em đã học}\par\smallskip
}{\par\medskip}
\newenvironment{emcothe}{%
  \par\medskip\noindent\textbf{Em có thể}\par\smallskip
}{\par\medskip}
\newenvironment{kienthuc}{%
  \par\medskip\noindent\textbf{Kiến thức cốt lõi}\par\smallskip
}{\par\medskip}
\newenvironment{traloi}{%
  \par\smallskip\noindent\textit{Trả lời / ghi chú của em}\par
  \vspace{1.2em}%
}{\par\medskip}
\makeatother

% từ gốc repo:
%   % Định nghĩa lệnh Lý thuyết — dùng khi biên dịch lt.tex bằng TeX.
% App web đọc lt.tex trực tiếp (bỏ \input file này).
% Trong lt.tex, từ thư mục bài:
%   \input{../../../../_lenh/lythuyet.tex}
% từ gốc repo:
%   \input{ngan-hang/_lenh/lythuyet.tex}
\makeatletter
\providecommand{\indam}[1]{\textbf{#1}}
\providecommand{\chuy}[1]{\par\smallskip\noindent\textit{#1}\par}
\providecommand{\luuy}[1]{\par\smallskip\noindent\textbf{Lưu ý.} #1\par}
\providecommand{\ghichu}[1]{\par\smallskip\noindent\textbf{Ghi chú.} #1\par}
\providecommand{\vidu}[1]{\par\smallskip\noindent\textbf{Ví dụ.} #1\par}
\providecommand{\Blythuyet}{\subsection*{Lý thuyết}}
% Hai cột: kiến thức | hình
\providecommand{\haicotchay}[2]{%
  \par\medskip\noindent
  \begin{minipage}[t]{0.52\linewidth}#1\end{minipage}\hfill
  \begin{minipage}[t]{0.46\linewidth}#2\end{minipage}%
  \par\medskip
}
\providecommand{\kienthuccot}[2]{\haicotchay{#1}{#2}}
% Dự phòng nếu chưa nạp graphicx / caption (không ghi đè bản thật)
\providecommand{\resizebox}[3]{#3}
\providecommand{\captionof}[2]{\par\smallskip\noindent\textit{#2}\par}
\newenvironment{hoatdong}[1]{%
  \par\medskip\noindent\textbf{Hoạt động.} #1\par\smallskip
}{\par\medskip}
\newenvironment{emcobiet}{%
  \par\medskip\noindent\textbf{Em có biết}\par\smallskip
}{\par\medskip}
\newenvironment{emdahoc}{%
  \par\medskip\noindent\textbf{Em đã học}\par\smallskip
}{\par\medskip}
\newenvironment{emcothe}{%
  \par\medskip\noindent\textbf{Em có thể}\par\smallskip
}{\par\medskip}
\newenvironment{kienthuc}{%
  \par\medskip\noindent\textbf{Kiến thức cốt lõi}\par\smallskip
}{\par\medskip}
\newenvironment{traloi}{%
  \par\smallskip\noindent\textit{Trả lời / ghi chú của em}\par
  \vspace{1.2em}%
}{\par\medskip}
\makeatother

\makeatletter
\providecommand{\indam}[1]{\textbf{#1}}
\providecommand{\chuy}[1]{\par\smallskip\noindent\textit{#1}\par}
\providecommand{\luuy}[1]{\par\smallskip\noindent\textbf{Lưu ý.} #1\par}
\providecommand{\ghichu}[1]{\par\smallskip\noindent\textbf{Ghi chú.} #1\par}
\providecommand{\vidu}[1]{\par\smallskip\noindent\textbf{Ví dụ.} #1\par}
\providecommand{\Blythuyet}{\subsection*{Lý thuyết}}
% Hai cột: kiến thức | hình
\providecommand{\haicotchay}[2]{%
  \par\medskip\noindent
  \begin{minipage}[t]{0.52\linewidth}#1\end{minipage}\hfill
  \begin{minipage}[t]{0.46\linewidth}#2\end{minipage}%
  \par\medskip
}
\providecommand{\kienthuccot}[2]{\haicotchay{#1}{#2}}
% Dự phòng nếu chưa nạp graphicx / caption (không ghi đè bản thật)
\providecommand{\resizebox}[3]{#3}
\providecommand{\captionof}[2]{\par\smallskip\noindent\textit{#2}\par}
\newenvironment{hoatdong}[1]{%
  \par\medskip\noindent\textbf{Hoạt động.} #1\par\smallskip
}{\par\medskip}
\newenvironment{emcobiet}{%
  \par\medskip\noindent\textbf{Em có biết}\par\smallskip
}{\par\medskip}
\newenvironment{emdahoc}{%
  \par\medskip\noindent\textbf{Em đã học}\par\smallskip
}{\par\medskip}
\newenvironment{emcothe}{%
  \par\medskip\noindent\textbf{Em có thể}\par\smallskip
}{\par\medskip}
\newenvironment{kienthuc}{%
  \par\medskip\noindent\textbf{Kiến thức cốt lõi}\par\smallskip
}{\par\medskip}
\newenvironment{traloi}{%
  \par\smallskip\noindent\textit{Trả lời / ghi chú của em}\par
  \vspace{1.2em}%
}{\par\medskip}
\makeatother

% từ gốc repo:
%   % Định nghĩa lệnh Lý thuyết — dùng khi biên dịch lt.tex bằng TeX.
% App web đọc lt.tex trực tiếp (bỏ \input file này).
% Trong lt.tex, từ thư mục bài:
%   % Định nghĩa lệnh Lý thuyết — dùng khi biên dịch lt.tex bằng TeX.
% App web đọc lt.tex trực tiếp (bỏ \input file này).
% Trong lt.tex, từ thư mục bài:
%   \input{../../../../_lenh/lythuyet.tex}
% từ gốc repo:
%   \input{ngan-hang/_lenh/lythuyet.tex}
\makeatletter
\providecommand{\indam}[1]{\textbf{#1}}
\providecommand{\chuy}[1]{\par\smallskip\noindent\textit{#1}\par}
\providecommand{\luuy}[1]{\par\smallskip\noindent\textbf{Lưu ý.} #1\par}
\providecommand{\ghichu}[1]{\par\smallskip\noindent\textbf{Ghi chú.} #1\par}
\providecommand{\vidu}[1]{\par\smallskip\noindent\textbf{Ví dụ.} #1\par}
\providecommand{\Blythuyet}{\subsection*{Lý thuyết}}
% Hai cột: kiến thức | hình
\providecommand{\haicotchay}[2]{%
  \par\medskip\noindent
  \begin{minipage}[t]{0.52\linewidth}#1\end{minipage}\hfill
  \begin{minipage}[t]{0.46\linewidth}#2\end{minipage}%
  \par\medskip
}
\providecommand{\kienthuccot}[2]{\haicotchay{#1}{#2}}
% Dự phòng nếu chưa nạp graphicx / caption (không ghi đè bản thật)
\providecommand{\resizebox}[3]{#3}
\providecommand{\captionof}[2]{\par\smallskip\noindent\textit{#2}\par}
\newenvironment{hoatdong}[1]{%
  \par\medskip\noindent\textbf{Hoạt động.} #1\par\smallskip
}{\par\medskip}
\newenvironment{emcobiet}{%
  \par\medskip\noindent\textbf{Em có biết}\par\smallskip
}{\par\medskip}
\newenvironment{emdahoc}{%
  \par\medskip\noindent\textbf{Em đã học}\par\smallskip
}{\par\medskip}
\newenvironment{emcothe}{%
  \par\medskip\noindent\textbf{Em có thể}\par\smallskip
}{\par\medskip}
\newenvironment{kienthuc}{%
  \par\medskip\noindent\textbf{Kiến thức cốt lõi}\par\smallskip
}{\par\medskip}
\newenvironment{traloi}{%
  \par\smallskip\noindent\textit{Trả lời / ghi chú của em}\par
  \vspace{1.2em}%
}{\par\medskip}
\makeatother

% từ gốc repo:
%   % Định nghĩa lệnh Lý thuyết — dùng khi biên dịch lt.tex bằng TeX.
% App web đọc lt.tex trực tiếp (bỏ \input file này).
% Trong lt.tex, từ thư mục bài:
%   \input{../../../../_lenh/lythuyet.tex}
% từ gốc repo:
%   \input{ngan-hang/_lenh/lythuyet.tex}
\makeatletter
\providecommand{\indam}[1]{\textbf{#1}}
\providecommand{\chuy}[1]{\par\smallskip\noindent\textit{#1}\par}
\providecommand{\luuy}[1]{\par\smallskip\noindent\textbf{Lưu ý.} #1\par}
\providecommand{\ghichu}[1]{\par\smallskip\noindent\textbf{Ghi chú.} #1\par}
\providecommand{\vidu}[1]{\par\smallskip\noindent\textbf{Ví dụ.} #1\par}
\providecommand{\Blythuyet}{\subsection*{Lý thuyết}}
% Hai cột: kiến thức | hình
\providecommand{\haicotchay}[2]{%
  \par\medskip\noindent
  \begin{minipage}[t]{0.52\linewidth}#1\end{minipage}\hfill
  \begin{minipage}[t]{0.46\linewidth}#2\end{minipage}%
  \par\medskip
}
\providecommand{\kienthuccot}[2]{\haicotchay{#1}{#2}}
% Dự phòng nếu chưa nạp graphicx / caption (không ghi đè bản thật)
\providecommand{\resizebox}[3]{#3}
\providecommand{\captionof}[2]{\par\smallskip\noindent\textit{#2}\par}
\newenvironment{hoatdong}[1]{%
  \par\medskip\noindent\textbf{Hoạt động.} #1\par\smallskip
}{\par\medskip}
\newenvironment{emcobiet}{%
  \par\medskip\noindent\textbf{Em có biết}\par\smallskip
}{\par\medskip}
\newenvironment{emdahoc}{%
  \par\medskip\noindent\textbf{Em đã học}\par\smallskip
}{\par\medskip}
\newenvironment{emcothe}{%
  \par\medskip\noindent\textbf{Em có thể}\par\smallskip
}{\par\medskip}
\newenvironment{kienthuc}{%
  \par\medskip\noindent\textbf{Kiến thức cốt lõi}\par\smallskip
}{\par\medskip}
\newenvironment{traloi}{%
  \par\smallskip\noindent\textit{Trả lời / ghi chú của em}\par
  \vspace{1.2em}%
}{\par\medskip}
\makeatother

\makeatletter
\providecommand{\indam}[1]{\textbf{#1}}
\providecommand{\chuy}[1]{\par\smallskip\noindent\textit{#1}\par}
\providecommand{\luuy}[1]{\par\smallskip\noindent\textbf{Lưu ý.} #1\par}
\providecommand{\ghichu}[1]{\par\smallskip\noindent\textbf{Ghi chú.} #1\par}
\providecommand{\vidu}[1]{\par\smallskip\noindent\textbf{Ví dụ.} #1\par}
\providecommand{\Blythuyet}{\subsection*{Lý thuyết}}
% Hai cột: kiến thức | hình
\providecommand{\haicotchay}[2]{%
  \par\medskip\noindent
  \begin{minipage}[t]{0.52\linewidth}#1\end{minipage}\hfill
  \begin{minipage}[t]{0.46\linewidth}#2\end{minipage}%
  \par\medskip
}
\providecommand{\kienthuccot}[2]{\haicotchay{#1}{#2}}
% Dự phòng nếu chưa nạp graphicx / caption (không ghi đè bản thật)
\providecommand{\resizebox}[3]{#3}
\providecommand{\captionof}[2]{\par\smallskip\noindent\textit{#2}\par}
\newenvironment{hoatdong}[1]{%
  \par\medskip\noindent\textbf{Hoạt động.} #1\par\smallskip
}{\par\medskip}
\newenvironment{emcobiet}{%
  \par\medskip\noindent\textbf{Em có biết}\par\smallskip
}{\par\medskip}
\newenvironment{emdahoc}{%
  \par\medskip\noindent\textbf{Em đã học}\par\smallskip
}{\par\medskip}
\newenvironment{emcothe}{%
  \par\medskip\noindent\textbf{Em có thể}\par\smallskip
}{\par\medskip}
\newenvironment{kienthuc}{%
  \par\medskip\noindent\textbf{Kiến thức cốt lõi}\par\smallskip
}{\par\medskip}
\newenvironment{traloi}{%
  \par\smallskip\noindent\textit{Trả lời / ghi chú của em}\par
  \vspace{1.2em}%
}{\par\medskip}
\makeatother

\makeatletter
\providecommand{\indam}[1]{\textbf{#1}}
\providecommand{\chuy}[1]{\par\smallskip\noindent\textit{#1}\par}
\providecommand{\luuy}[1]{\par\smallskip\noindent\textbf{Lưu ý.} #1\par}
\providecommand{\ghichu}[1]{\par\smallskip\noindent\textbf{Ghi chú.} #1\par}
\providecommand{\vidu}[1]{\par\smallskip\noindent\textbf{Ví dụ.} #1\par}
\providecommand{\Blythuyet}{\subsection*{Lý thuyết}}
% Hai cột: kiến thức | hình
\providecommand{\haicotchay}[2]{%
  \par\medskip\noindent
  \begin{minipage}[t]{0.52\linewidth}#1\end{minipage}\hfill
  \begin{minipage}[t]{0.46\linewidth}#2\end{minipage}%
  \par\medskip
}
\providecommand{\kienthuccot}[2]{\haicotchay{#1}{#2}}
% Dự phòng nếu chưa nạp graphicx / caption (không ghi đè bản thật)
\providecommand{\resizebox}[3]{#3}
\providecommand{\captionof}[2]{\par\smallskip\noindent\textit{#2}\par}
\newenvironment{hoatdong}[1]{%
  \par\medskip\noindent\textbf{Hoạt động.} #1\par\smallskip
}{\par\medskip}
\newenvironment{emcobiet}{%
  \par\medskip\noindent\textbf{Em có biết}\par\smallskip
}{\par\medskip}
\newenvironment{emdahoc}{%
  \par\medskip\noindent\textbf{Em đã học}\par\smallskip
}{\par\medskip}
\newenvironment{emcothe}{%
  \par\medskip\noindent\textbf{Em có thể}\par\smallskip
}{\par\medskip}
\newenvironment{kienthuc}{%
  \par\medskip\noindent\textbf{Kiến thức cốt lõi}\par\smallskip
}{\par\medskip}
\newenvironment{traloi}{%
  \par\smallskip\noindent\textit{Trả lời / ghi chú của em}\par
  \vspace{1.2em}%
}{\par\medskip}
\makeatother

\subsubsection{Đối tượng nghiên cứu của Vật lí và mục tiêu của môn Vật lí}
\paragraph{Đối tượng nghiên cứu của Vật lí}
\begin{itemize}
\item Vật lí là ngành khoa học nghiên cứu về các dạng vận động của vật chất và năng lượng ở mọi cấp độ khác nhau.
\item Ở cấp độ vi mô, Vật lí tập trung nghiên cứu thế giới của các hạt cấu trúc siêu nhỏ như nguyên tử, phân tử, hạt nhân và các hạt cơ bản.
\item Ở cấp độ vĩ mô, Vật lí nghiên cứu thế giới của các vật thể thông thường xung quanh chúng ta, Trái Đất, các hành tinh, Mặt Trời, Hệ Mặt Trời, các thiên hà, siêu thiên hà và toàn bộ Vũ trụ rộng lớn.
\end{itemize}
\paragraph{Mục tiêu của môn Vật lí}
\begin{itemize}
\item Mục tiêu cơ bản của Vật lí là khám phá ra quy luật tổng quát nhất chi phối sự vận động của vật chất và năng lượng, cũng như các tương tác giữa chúng ở mọi cấp độ từ vi mô đến vĩ mô.
\item Học tập môn Vật lí giúp học sinh hình thành, phát triển năng lực vật lí với các biểu hiện chính bao gồm: nhận thức vật lí, tìm hiểu tự nhiên dưới góc độ vật lí, và vận dụng kiến thức, kĩ năng đã học vào thực tiễn cuộc sống.
\end{itemize}

\subsubsection{Vật lí với cuộc sống, khoa học, kĩ thuật và công nghệ}
\paragraph{Vật lí với cuộc sống}
\begin{itemize}
\item Tri thức vật lí có ảnh hưởng vô cùng rộng lớn, là cơ sở khoa học để con người chế tạo và giải thích nguyên tắc hoạt động của vô số thiết bị, vật dụng tiện ích hằng ngày.
\item Các thiết bị quen thuộc trong gia đình như nồi cơm điện, bếp từ, lò vi sóng, tủ lạnh, điều hòa nhiệt độ, tivi, điện thoại thông minh đều hoạt động dựa trên các nguyên lí vật lí.
\end{itemize}
\paragraph{Vật lí với khoa học, kĩ thuật và công nghệ}
\begin{itemize}
\item Vật lí được coi là cơ sở, là nền tảng của khoa học tự nhiên và công nghệ hiện đại. Các kiến thức vật lí được áp dụng rộng rãi để thúc đẩy sự ra đời và phát triển của các ngành công nghiệp, y học, nông nghiệp và truyền thông.
\item Lịch sử loài người đã trải qua bốn cuộc cách mạng công nghiệp lớn, mỗi cuộc cách mạng đều gắn liền và được thúc đẩy bởi những phát minh mang tính đột phá của Vật lí:
\begin{itemize}
\item Cuộc cách mạng công nghiệp lần thứ nhất (thế kỉ XVIII): Xuất hiện các thiết bị máy móc nhờ việc phát minh ra động cơ hơi nước, dựa trên những nghiên cứu sâu sắc về Nhiệt học và Nhiệt động lực học.
\item Cuộc cách mạng công nghiệp lần thứ hai (thế kỉ XIX): Đặc trưng bởi sự xuất hiện của các thiết bị dùng điện trong mọi lĩnh vực sản xuất và đời sống, dựa trên các khám phá về Điện từ học.
\item Cuộc cách mạng công nghiệp lần thứ ba (thế kỉ XX): Gắn liền với sự ra đời của máy tính, internet, công nghệ bán dẫn và tự động hóa sản xuất, dựa trên các thành tựu của Vật lí lượng tử và Vật lí chất rắn.
\item Cuộc cách mạng công nghiệp lần thứ tư (thế kỉ XXI): Là sự phát triển như vũ bão của trí tuệ nhân tạo (AI), robot thông minh, internet vạn vật (IoT), công nghệ sinh học và công nghệ vật liệu siêu nhỏ, siêu nhẹ (nano), dựa trên các nghiên cứu tiên tiến của Vật lí hiện đại.
\end{itemize}
\end{itemize}
\begin{emcobiet}
Việc ứng dụng các thành tựu của vật lí vào công nghệ không chỉ mang lại những lợi ích to lớn cho nhân loại mà còn có thể gây ra những tác động tiêu cực như ô nhiễm môi trường sống, hủy hoại hệ sinh thái và biến đổi khí hậu nếu con người không sử dụng đúng phương pháp, đúng mục đích hoặc lạm dụng chúng cho các hoạt động hủy diệt.
\end{emcobiet}

\subsubsection{Phương pháp nghiên cứu Vật lí}
\paragraph{Phương pháp thực nghiệm}
\begin{itemize}
\item Phương pháp thực nghiệm là phương pháp sử dụng các thí nghiệm khoa học để phát hiện kết quả mới nhằm kiểm chứng, hoàn thiện, bổ sung hoặc bác bỏ một giả thuyết nào đó.
\item Kết quả mới phát hiện từ thực nghiệm cần được giải thích bằng các lí thuyết vật lí đã biết hoặc xây dựng một lí thuyết mới phù hợp.
\item Đây là phương pháp có tính chất quyết định đối với sự hình thành và phát triển của Vật lí học.
\item Điển hình là cuộc cách mạng nghiên cứu thực nghiệm của Galileo Galilei khi thả rơi tự do hai quả cầu có khối lượng khác nhau tại tháp nghiêng Pisa để chứng minh chúng rơi chạm đất cùng lúc, bác bỏ quan niệm suy luận chủ quan tồn tại hơn 2000 năm của Aristotle.
\end{itemize}
\paragraph{Phương pháp mô hình}
\begin{itemize}
\item Phương pháp mô hình sử dụng các mô hình đơn giản hóa để thay thế cho vật thật nhằm nghiên cứu, giải thích các tính chất của vật thật và tìm ra cơ chế hoạt động của nó.
\item Các loại mô hình thường dùng trong Vật lí bao gồm:
\begin{itemize}
\item Mô hình vật chất: Là các vật thu nhỏ hoặc phóng to của vật thật, ví dụ quả địa cầu là mô hình thu nhỏ của Trái Đất, mô hình cấu trúc tinh thể, mô hình mẫu hành tinh nguyên tử của Rutherford.
\item Mô hình lí thuyết: Khi nghiên cứu chuyển động của một ô tô chạy đường dài trên quãng đường hàng trăm kilômét, người ta coi ô tô là một "chất điểm"; hoặc để biểu diễn đường truyền của ánh sáng, người ta dùng mô hình "tia sáng".
\item Mô hình toán học: Sử dụng các công thức, phương trình, đồ thị để mô tả các đặc điểm của đối tượng nghiên cứu, ví dụ công thức tính quãng đường chuyển động thẳng đều $s = v \cdot t$.
\end{itemize}
\end{itemize}
\begin{hoatdong}{Sơ đồ hóa tiến trình tìm hiểu thế giới tự nhiên dưới góc độ vật lí}
Hãy thực hiện thảo luận để sắp xếp đúng tiến trình nghiên cứu thế giới tự nhiên gồm 5 bước sau:
\begin{itemize}
\item Bước 1: Quan sát hiện tượng, suy luận để xác định đối tượng cần nghiên cứu và đề xuất vấn đề.
\item Bước 2: Hình thành giả thuyết khoa học về vấn đề cần nghiên cứu.
\item Bước 3: Thiết kế, xây dựng mô hình lí thuyết hoặc mô hình thực nghiệm để kiểm chứng giả thuyết.
\item Bước 4: Thực hiện tính toán theo mô hình lí thuyết hoặc tiến hành thí nghiệm thực nghiệm để thu thập, xử lí và phân tích số liệu nhằm kiểm chứng, điều chỉnh hoặc bác bỏ giả thuyết ban đầu.
\item Bước 5: Rút ra kết luận khoa học và đánh giá mô hình nghiên cứu.
\end{itemize}
\end{hoatdong}
\begin{emdahoc}
\begin{itemize}
\item Đối tượng nghiên cứu chủ yếu của Vật lí là các dạng vận động của vật chất và năng lượng ở mọi cấp độ khác nhau.
\item Vật lí học là cơ sở của khoa học tự nhiên, công nghệ và kĩ thuật hiện đại, có ảnh hưởng to lớn đến cuộc sống và sự phát triển của xã hội loài người.
\item Phương pháp nghiên cứu đặc trưng của Vật lí gồm hai phương pháp chính là phương pháp thực nghiệm và phương pháp lí thuyết, chúng luôn bổ trợ chặt chẽ cho nhau, trong đó phương pháp thực nghiệm mang tính quyết định.
\item Tiến trình nghiên cứu của Vật lí tuân theo các bước khoa học nghiêm ngặt: Quan sát - Đặt vấn đề - Giả thuyết - Kiểm chứng (bằng lí thuyết hoặc thực nghiệm) - Kết luận.
\end{itemize}
\end{emdahoc}
\begin{emcothe}
\begin{itemize}
\item Nhận diện và liên hệ được các hiện tượng tự nhiên xung quanh với các quy luật động học, nhiệt học, điện học và quang học đã học.
\item Thảo luận và đánh giá được vai trò của Vật lí trong việc hình thành các công nghệ xanh, bảo vệ môi trường sống và ứng phó với biến đổi khí hậu trong thế kỉ mới.
\end{itemize}
\end{emcothe}